\documentclass[10pt,a4paper]{article}
\usepackage{sectsty}
\usepackage{hyperref}
\sectionfont{\raggedright}

\newcommand{\guidance}[1]{\par\smallskip\noindent\textit{#1}\par\smallskip}

\title{Written Report Template}
\author{
  Alessandro ARENSBERG\\
  \texttt{aarensberg@albertschool.com}%
  \and
  Charles DEFOSSEUX\\
  \texttt{cdefosseux@albertschool.com}%
  \and
  Louis ALLIO\\
  \texttt{lallio@albertschool.com}
}
\date{}

\begin{document}
\maketitle

\begin{abstract}
\guidance{Summarize the research question, data and graph model, analytical methods, main findings, and conclusion.}
\end{abstract}

\section{Introduction}
\guidance{Introduce the topic and motivation, state the objective, and explain its relevance.}
\noindent Example citation: Newman surveys complex networks \cite{newman2003}. In the source, this is written as \texttt{\textbackslash cite\{newman2003\}}.


\section{Dataset Presentation}
\guidance{Describe the dataset and its source, explain how it is represented as a graph, and report relevant graph statistics, including the numbers of nodes and edges.}


\section{Description of the Analysis}
\guidance{Explain the methods, software, and parameter choices; describe how you assess connectivity, distances, centrality, or communities as relevant. Disclose any reused code and whether AI tools were used, stating their purpose.}


\section{Results}
\guidance{Present the results clearly, using tables or figures where useful, and explain what they show in relation to your analysis.}


\section{Conclusions}
\guidance{Summarize the main findings and note relevant limitations or possible future work.}


\guidance{Add references to \texttt{biblio.bib} and cite them in the text using \texttt{\textbackslash cite\{key\}}. Include sources for data, software, methods, and reused code.}
\bibliographystyle{plain}
\bibliography{biblio}


\end{document}